% !PRINTED 625
\weberpage{625}
%<<CONT
shall vanish then and only then, when the function $F(\theta,z)$ occurs among its irreducible factors. This is the definition of vanishing, according to which, if one wished to be entirely correct, every equation $G(\theta,z) = 0$ would have to be understood as a congruence with respect to the modulus $F$. In order not to be prolix, however, we shall not employ this mode of expression here.

If we divide equation (1) by $a_0$, it takes the form
\weq{3}{f(\theta,z) = \theta^n + b_1\theta^{n-1} + b_2\theta^{n-2} + \cdots + b_{n-1}\theta + b_n,}
in which the $b_1, b_2, \ldots, b_n$ are integral or fractional rational functions of $z$.

The system of all integral and fractional rational functions $\Phi(\theta,z)$ of $\theta$ and $z$ in which the denominator is not divisible by $F$, and so does not vanish, has the property of reproducing itself by addition, subtraction, multiplication and division (except by 0), and therefore forms a \gesperrt{field of algebraic functions} (Bd.\,I, \S\,146), which I shall denote by $\Omega$.

The degree $n$ of the irreducible equation (2) or (3), that is, of the equation of lowest degree which $\theta$ satisfies, is called the \gesperrt{degree of the algebraic field with respect to} $\theta$, or also the \gesperrt{degree of} $\theta$.

If $\varphi(\theta)$ is an integral function of $\theta$ whose coefficients are integral or fractional rational functions of $z$, then by division one can form an equation:
\[
\varphi(\theta) = f(\theta)q(\theta) + r(\theta),
\]
in which $q(\theta)$ and $r(\theta)$ are functions of the same kind as $\varphi$, but of which $r(\theta)$ is \gesperrt{at most of degree} $n-1$. By reason of (3) we then have
\weq{4}{\varphi(\theta) = r(\theta).}

If $\varphi(\theta)$ is not divisible by $f(\theta)$, then the two functions have no rational divisor in common, and one can determine two functions of the field $f_1(\theta), \varphi_1(\theta)$ in such a way that
\[
f(\theta)f_1(\theta) + \varphi(\theta)\varphi_1(\theta) = 1
\]
holds (Bd.\,I, \S\,6). Since now $f(\theta) = 0$, there follows from this:
\weq{5}{\varphi_1(\theta) = \frac{1}{\varphi(\theta)}.}
